\documentclass[11pt,a4paper]{article}
\usepackage[margin=2.4cm]{geometry}
\usepackage{amsmath,amssymb}
\usepackage{enumitem}
\usepackage{xcolor}
\usepackage{hyperref}
\hypersetup{colorlinks=true,linkcolor=blue!55!black,urlcolor=blue!55!black}
\newcommand{\code}[1]{\texttt{#1}}
\newcommand{\uvec}{\mathbf{u}}
\newcommand{\zhat}{\hat{\mathbf{z}}}
\setlength{\parskip}{0.35em}
\setlength{\parindent}{0pt}

\title{\vspace{-1.5em}\textbf{A balance-consistent restart for the moving-cavity mesh change}\\[0.3em]
\large Deriving the conditions newly-opened and newly-iced cells must satisfy\\
for settle-free continuation}
\author{AWI-ESM3\,$\leftrightarrow$\,PISM moving cavity --- design note}
\date{2026-07-07}

\begin{document}
\maketitle

\begin{abstract}
\noindent
On every moving-cavity leg the FESOM restart is remapped onto a sub-mesh whose ice front has
moved. Empirically the leg then blows up within days: \emph{fast} as a vertical-CFL runaway at
$\Delta t=1200$\,s, or \emph{slow} as a barotropic $\eta_n\!\to\!$NaN at $\Delta t=60$\,s. We have
falsified every \emph{value-only} restart fix (removing static density inversions; nearest-neighbour
velocity; running without a settle). This note shows that the two blow-ups are the two faces of a
single \textbf{geostrophic adjustment}, that the balanced state a stable restart must already be in
is \textbf{fully determined} by the density and sea-surface height together with the surrounding
(already-spun-up) water masses, and that we can therefore \textbf{compute} the required velocity
(thermal wind $+$ barotropic reference) rather than guess it. The corollary is that the remap must
also not over-sharpen the opened-column fronts. The plan is to build both into the F90 remap and
verify the balance metrics \emph{before} the run, giving settle-free continuation.
\end{abstract}

\section{Status: what is falsified}
The instability is a coastal, mid-column ($\sim$95\,m) mode that grows over days from the changed
margin. Ruled out, each verified at the data level and each leaving the blow-up essentially
unchanged:
\begin{itemize}[leftmargin=1.5em,itemsep=1pt]
  \item \textbf{Static density inversions} --- a coherent-column $T/S$ refill drove the worst
        stability over all retreated cells from $-8.8$ to $-0.5$\,kg\,m$^{-3}$ (11 severe $\to$ 0);
        the leg still blew (\code{mstep}$\approx$164, CFLz$\approx$8.5).
  \item \textbf{Nearest-neighbour velocity} (Finn Heukamp / SO-ASE handling) --- filling changed
        elements with the nearest old element's current (no cell left at zero) lowered peak CFLz
        $8.5\to5.7$ but \emph{still} blew at \code{mstep}$\approx$164, now with a barotropic
        $\eta_n\!\to\!$NaN as well.
  \item \textbf{No settle / timestep} --- $\Delta t=60$\,s removed the fast CFLz mode entirely
        (0 warnings, $\sim$10 days) but a slow barotropic $\eta_n$ mode grew and NaN'd.
\end{itemize}
A value-only restart cannot be made \emph{balanced}: the imbalance is lateral (fronts) $+$ missing
momentum, not a vertical inversion. We now derive what "balanced" means quantitatively.

\section{The balance a stable restart must satisfy}
Started from a density field, the primitive equations relax to geostrophic--hydrostatic balance;
"no adjustment shock" is equivalent to the restart already satisfying, at every wet point:
\begin{align}
  \text{(1) static stability:}\quad & N^2 = -\tfrac{g}{\rho_0}\,\partial_z\rho \;\ge\; 0, \\
  \text{(2) hydrostatic:}\quad & p(z) = \rho_0 g\,\eta \;+\; \underbrace{\rho_i g\,h_i}_{\text{ice-shelf load}} \;+\; \int_z^{0}\!\rho\,g\,dz', \\
  \text{(3) geostrophy:}\quad & f\,\zhat\times\uvec = -\tfrac{1}{\rho_0}\nabla_h p, \\
  \text{(4) thermal wind:}\quad & \partial_z\uvec = -\tfrac{g}{f\rho_0}\,\zhat\times\nabla_h\rho, \\
  \text{(5) barotropic reference:}\quad & \uvec(z{=}0) = \tfrac{g}{f}\,\zhat\times\nabla_h\eta, \\
  \text{(6) non-divergence:}\quad & \nabla_h\cdot\uvec \approx 0 .
\end{align}
(1) and (6) are the two we \emph{watch}: (1) keeps convection off, and (6) --- automatically
satisfied by geostrophic flow to $O(\text{Rossby})$ --- keeps the diagnosed $w$ (hence CFLz)
\emph{and} $\partial_t\eta$ (hence the barotropic NaN) small. So a state that satisfies (3)--(5)
kills \emph{both} the fast and the slow blow-up at once.

\paragraph{The velocity is computable, not guessable.} Integrating (4) from the surface and adding
the barotropic reference (5),
\begin{equation}
  \boxed{\;\uvec(z) \;=\; \underbrace{\tfrac{g}{f}\,\zhat\times\nabla_h\eta}_{\text{barotropic, from }\eta}
  \;+\; \underbrace{\tfrac{g}{f\rho_0}\int_z^{0}\zhat\times\nabla_h\rho\,\,dz'}_{\text{thermal wind, from }\rho(T,S)}\;}
  \label{eq:ug}
\end{equation}
with the vertical integration constant / reference fixed by the \textbf{surrounding unchanged
cells}, which are already in the model's balanced state. Set $T,S,\eta$ from the neighbours; then
$\uvec$ is \emph{determined} --- it must not be zeroed (v4) nor NN-copied (v5).

\section{Quantitative diagnosis: the missing jet}
At the Weddell seed the fill leaves a front of $\Delta S\approx3$\,psu over $\Delta x\approx25$\,km,
so $\Delta\rho\approx\beta_S\Delta S\approx2.4$\,kg\,m$^{-3}$ and $\nabla_h\rho\approx1.0\times10^{-4}$\,kg\,m$^{-4}$.
With $f(-66^\circ)\approx-1.33\times10^{-4}$\,s$^{-1}$, (4) gives
\begin{equation}
  \partial_z u = \tfrac{g}{f\rho_0}\,\nabla_h\rho \approx 7\times10^{-3}\ \text{s}^{-1}
  \;\Longrightarrow\; \Delta u \approx 0.7\ \text{m\,s}^{-1}\ \text{over }100\text{ m.}
\end{equation}
\textbf{The balanced front carries a $\sim$0.7\,m\,s$^{-1}$ geostrophic jet.} We initialised it at
$0$ (v4) or $\sim$0.06\,m\,s$^{-1}$ (v5): an $O(0.7\,\text{m s}^{-1})$ momentum deficit --- precisely
the violent adjustment observed. NN cannot supply it because the required current is specific to
\emph{this} density gradient, not the neighbour's background flow.

\section{Conditions per case}
\paragraph{Added points (ice retreated / newly wet).} At each newly-wet level:
\begin{enumerate}[leftmargin=1.6em,itemsep=1pt]
  \item $T,S$: a statically stable ($N^2\!\ge\!0$) blend of the surrounding water masses at that
        depth --- no grid-scale front sharper than the physical one.
  \item $\uvec$: the thermal-wind profile of Eq.~\eqref{eq:ug} implied by that $T/S$, referenced to
        the neighbouring balanced current.
  \item $\eta$: smooth to the neighbours (hydrostatically consistent).
  \item \code{hnode}: ALE-consistent (sums to $H+\eta$), so no spurious $\partial_t h\to w$.
\end{enumerate}
\paragraph{Points next to removed points (ice advanced / newly iced).} The new ice adds a surface
load $\rho_i g h_i$, so $\eta$ at the iced cell drops hydrostatically by
$\approx(\rho_i/\rho_0)h_i$. The \emph{adjacent} open cells must see that pressure step as a
\emph{balanced} geostrophic current: their $\uvec$ must include the geostrophic response
(5) to the new ice-loading slope. Computable from the PISM draft $h_i$ and hydrostaty.

\section{Corollary: do not over-sharpen the front}
A $0.7$\,m\,s$^{-1}$ jet is high for an Antarctic cavity front, which flags that the current fill
\textbf{over-sharpens} it: the shelf-base "hold upward" plants $S\!\approx\!31$ through the whole
newly-opened column, where physically only a thin near-base melt layer is that fresh. The two levers
are therefore linked:
\[
  \text{gentler }T/S\text{ front}\ \Rightarrow\ \text{smaller }\nabla_h\rho\ \Rightarrow\
  \text{smaller required jet;}\qquad
  \text{set that jet by thermal wind}\ \Rightarrow\ \text{the remainder is in balance.}
\]
Do both and conditions (1)--(6) hold at $t=0$ by construction $\Rightarrow$ no adjustment
$\Rightarrow$ \textbf{settle-free continuation}, i.e.\ the full year to the next PISM call at
$1200$\,s.

\section{Implementation plan (F90 remap, \code{mo\_remap\_fields.F90})}
\begin{enumerate}[leftmargin=1.6em,itemsep=2pt]
  \item \textbf{Soften the opened-column $T/S$.} Replace the shelf-base "hold upward" with a
        distance/depth-weighted blend toward the surrounding open-ocean profile (keep real cavity
        water near the base; relax the mid/upper column), preserving $N^2\!\ge\!0$.
  \item \textbf{Density gradient.} On the new sub-mesh compute $\nabla_h\rho$ per level from the
        (softened) $T,S$ via the finite-volume gradient FESOM already uses (nabla of a nodal
        field on the triangulation); use a linearised EOS locally (sufficient for the gradient).
  \item \textbf{Thermal-wind integral.} Integrate Eq.~\eqref{eq:ug} downward from the surface on
        each changed \emph{element}, with the integration reference set to the nearest unchanged
        element's velocity (the SO-ASE NN, now used only as the \emph{barotropic reference}, not
        the whole field).
  \item \textbf{Barotropic part.} Add $\tfrac{g}{f}\zhat\times\nabla_h\eta$ from the (donor-filled,
        ice-loaded) $\eta$; near-coast $f$ is safely large (Antarctic), no equatorial issue.
  \item \textbf{Verify \emph{before} running.} Report, over the changed cells: $\min N^2$,
        $\max|\nabla_h\cdot\uvec|$ (target $\ll$ the CFL bound), and the thermal-wind residual
        $|\partial_z\uvec + \tfrac{g}{f\rho_0}\zhat\times\nabla_h\rho|$. These are cheap, and a green
        check here predicts a stable leg.
  \item \textbf{Re-test} the movcav8 chunk-3 leg at $1200$\,s, no settle.
\end{enumerate}

\section{Caveats / open points}
\begin{itemize}[leftmargin=1.5em,itemsep=1pt]
  \item \textbf{Discrete divergence.} The mesh-discrete geostrophic velocity has a small residual
        $\nabla_h\cdot\uvec$; it should sit below the CFL bound but must be checked (step 5).
  \item \textbf{Barotropic--baroclinic consistency.} $\eta$ and the density's steric contribution
        must be mutually consistent; if the donor $\eta$ and the thermal-wind interior disagree at
        the reference level, prefer the neighbour-matched value and let $\eta$ follow.
  \item \textbf{Ageostrophic part.} Tidal/Ekman components are not in (3)--(5); at $t=0$ for a
        spun-up interior geostrophy dominates, and matching the neighbour current as the reference
        carries the ageostrophic background across.
  \item \textbf{Still-sub-ice retreat cells} (e.g.\ ulev $14\!\to\!2$, the actual seed) are inside
        the cavity, not open ocean; their $\uvec$ reference is the neighbouring \emph{cavity} flow,
        and their $T/S$ blend is toward cavity, not open-ocean, water.
  \item \textbf{Source term.} If the per-leg front motion is itself implausibly large (deep-draft
        open$\leftrightarrow$cavity transitions in one 10-yr PISM step), gentling the
        \code{cavity\_nlvls} regeneration reduces $\nabla_h\rho$ at the root and makes every step
        above easier.
\end{itemize}

\end{document}
